\section{Indirect Rebound Effects}
\label{sec:indirect}

Indirect rebound effects arise when AI-driven efficiency gains
in one sector stimulate increased resource consumption in other
areas of the economy. These effects propagate through changes in
behavior, industrial structure, and economic incentives, and may
exceed direct effects in scale \cite{luccioni2025}.

A prominent example is AI-based logistics optimization. Systems
such as ORION, deployed by UPS, have saved over ten million
gallons of fuel annually \cite{mhlanga2025}. However, the
resulting cost reductions tend to stimulate higher order volumes
rather than reduce overall transport activity, potentially
offsetting the savings achieved \cite{mhlanga2025,luccioni2025}.

A further mechanism operates through AI-driven advertising and
recommendation systems. By increasing the precision and reach
of commercial messaging, these systems may contribute to
additional consumption that is rarely accounted for in
sustainability assessments \cite{luccioni2025}. A direct causal
link to aggregate consumption increases is difficult to isolate
empirically and should be treated as a plausible rather than
an established mechanism.

Empirical evidence from the Chinese digital economy suggests
that the relationship between digitalization and energy
consumption is nonlinear: below a critical threshold, efficiency
gains reduce energy use; once exceeded, rebound effects become
dominant, producing a net increase \cite{chen2025}. Research
on the broader ICT sector concludes that efficiency improvements
fail to account for rebound effects that can offset emission
savings or increase emissions overall \cite{widdicks2023}.

Indirect rebound effects are harder to quantify, manifesting
across sectors and over longer time horizons. The rebound
literature suggests that economy-wide effects may dominate the
direct channel \cite{sorrell2009}, implying that the cases
reviewed here represent a lower bound on overall rebound
magnitude.
